%% toc: P3 39 39.10
\section{Proofs: computable presentations of countable structures}
\label{sec:annex-p03-w03-p03-063-computable-presentations}

Throughout, \(\Lang\) is a computable signature in the sense of
\cref{def:p03-w03-p03-063-computable-signature}, and codes are those of
\cref{def:p03-w03-p03-063-sentence-coding}.  Two facts about the numerical
pairing map of \cref{sec:turing-machines-01} are used repeatedly.

\emph{(P1) Bijectivity.}  For each \(s\geq0\), the pairs \((a,b)\) with
\(a+b=s\) receive the codes \(s(s+1)/2+b\) for \(b=0,\ldots,s\).  These are
the consecutive integers from \(s(s+1)/2\) to \((s+1)(s+2)/2-1\).  The blocks
for \(s=0,1,2,\ldots\) partition \(\N\), so
\(\langle\cdot,\cdot\rangle_{\N}:\N^2\to\N\) is a bijection.  Its inverse is
computed by the decoding procedure in \cref{sec:turing-machines-01}.  By induction on \(k\),
\(\operatorname{code}_k:\N^k\to\N\) is a bijection for every \(k\geq1\), and
it is computably invertible; \(\operatorname{code}_0\) sends the empty tuple
to \(0\).

\emph{(P2) Size.}  Write \(T=(a+b)(a+b+1)/2\).  Then
\(\langle a,b\rangle_{\N}=T+b\geq b\).  If \(a+b\geq1\), then
\(T\geq a+b\geq a\); if \(a+b=0\), then \(a=0\).  Hence
\(\langle a,b\rangle_{\N}\geq a\).  If \(a\geq1\), then \(T\geq1\) and so
\(\langle a,b\rangle_{\N}>b\).  By induction on \(k\), every entry of a tuple
is at most its code \(\operatorname{code}_k\).

A function defined by course-of-values recursion, in which the value at \(n\)
is computed by a computable step from \(n\) and the values at arguments
smaller than \(n\), is computable.  The procedure computes the finite table of
values at \(0,1,\ldots,n\) in order, and each entry is obtained by a
terminating computation from the earlier ones.

\begin{annexitem}{annex:p03-w03-p03-063-parsing}{thm:p03-w03-p03-063-parsing}{Decidable syntax}
\begin{proof}
\emph{Smaller subcodes.}  In a code
\(\langle t,\langle i,\operatorname{code}_k(\vec z)\rangle_{\N}\rangle_{\N}\)
with tag \(t\in\{2,4\}\), (P2) gives
\(z_l\leq\operatorname{code}_k(\vec z)\leq\langle i,\operatorname{code}_k(\vec z)\rangle_{\N}\).
Since \(t\geq1\), the last number is strictly less than the whole code.  The
same argument applies to tag \(3\), with \(\langle\ulcorner t\urcorner,\ulcorner s\urcorner\rangle_{\N}\),
and to tag \(5\).  Hence every immediate subexpression has a strictly smaller
code.

\emph{Terms.}  The index sets \(I_R,I_F,I_C\) are initial segments of \(\N\),
so each is either \(\N\) or a finite set \(\{0,\ldots,p-1\}\), and each is
decidable.  Define \(\chi_{\mathrm{Tm}}(n)\) as follows.  Decode
\(n=\langle a,b\rangle_{\N}\) by (P1).  Put \(\chi_{\mathrm{Tm}}(n)=1\) iff
one of the following holds:
\begin{itemize}
\item \(a=0\);
\item \(a=1\) and \(b\in I_C\); or
\item \(a=2\), \(b=\langle i,z\rangle_{\N}\) with \(i\in I_F\), and, for
  \(k=a_F(i)\), either \(k=0\) and \(z=0\), or \(k\geq1\),
  \(z=\operatorname{code}_k(z_1,\ldots,z_k)\) and
  \(\chi_{\mathrm{Tm}}(z_l)=1\) for every \(l\).
\end{itemize}
The step is computable, because \(a_F\) is computable and decoding is
computable, and it consults only the smaller arguments \(z_l\).  So
\(\chi_{\mathrm{Tm}}\) is computable by course-of-values recursion.  An
induction on terms shows that \(\chi_{\mathrm{Tm}}(n)=1\) exactly when \(n\)
is the code of a term.

\emph{Atomic formulas and literals.}  Put \(n\in\mathrm{At}\) iff
\(n=\langle3,\langle u,v\rangle_{\N}\rangle_{\N}\) with \(u,v\in\mathrm{Tm}\),
or \(n=\langle4,\langle i,z\rangle_{\N}\rangle_{\N}\) with \(i\in I_R\) and
\(z\) the code of an \(a_R(i)\)-tuple of term codes.  Put
\(\mathrm{Lit}=\mathrm{At}\cup\{\langle5,m\rangle_{\N}:m\in\mathrm{At}\}\).
Both conditions are decidable from \(\chi_{\mathrm{Tm}}\), \(a_R\) and
decoding.

\emph{Unique readability.}  The tag determines the form of the expression.
The pairing map is injective by (P1), so the tag and its argument are
determined.  The index \(i\) then determines the arity \(k\), and
\(\operatorname{code}_k\) is injective.  Induction on codes now shows that
each code in \(\mathrm{Tm}\cup\mathrm{Lit}\) comes from exactly one
expression.

\emph{Occurrence.}  Define \(\mathrm{Occ}(j,n)\) by recursion on \(n\).  For
tag \(0\), it holds iff \(b=j\).  For tag \(1\), it is false.  For the tags
\(2\), \(3\), \(4\) and \(5\), it holds iff it holds for some immediate
subcode.  This is a course-of-values recursion in \(n\) with \(j\) as a
parameter, so the relation is decidable.  Since
\(\ulcorner x_j\urcorner=\langle0,j\rangle_{\N}\geq j\), and every subcode is
at most \(n\), every \(j\) with \(\mathrm{Occ}(j,n)\) satisfies \(j\leq n\).
Therefore the finite set of variable indices of \(n\) is computable from
\(n\) by checking \(j=0,\ldots,n\).
\end{proof}
\end{annexitem}

\begin{annexitem}{annex:p03-w03-p03-063-equivalence}{thm:p03-w03-p03-063-equivalence}{Three forms of a computable presentation}
\begin{proof}
Let \(P\) be the set of codes of true atomic \(\Lang(M)\)-sentences.  For a
literal code \(n\) whose variable indices lie in \(M\), let \(\sigma_n\) be
the \(\Lang(M)\)-sentence it codes.

\emph{(a) implies (b).}  If \(j\in M\), then \(c_j=c_j\) is a true atomic
\(\Lang(M)\)-sentence, so
\(\ulcorner x_j=x_j\urcorner\in\ulcorner\operatorname{Diag}(\mathcal M)\urcorner\).
If \(j\notin M\), then no code mentioning \(x_j\) belongs to that set.  Hence
\(j\in M\) iff
\(\ulcorner x_j=x_j\urcorner\in\ulcorner\operatorname{Diag}(\mathcal M)\urcorner\),
and the map \(j\mapsto\ulcorner x_j=x_j\urcorner\) is computable.  So \(M\)
is decidable.  Negated codes carry tag \(5\) and lie outside \(\mathrm{At}\),
so \(P=\ulcorner\operatorname{Diag}(\mathcal M)\urcorner\cap\mathrm{At}\).
This is decidable by \cref{thm:p03-w03-p03-063-parsing}.

\emph{(b) implies (c).}  Define
\(\rho(i,z)\) as follows.  If \(i\notin I_R\), return \(0\).  Otherwise put
\(k=a_R(i)\) and decode \(z=\operatorname{code}_k(n_1,\ldots,n_k)\) when
\(k\geq1\).  If some \(n_l\notin M\), return \(0\).  Otherwise return
\(\chi_P(\ulcorner\mathsf R_i(x_{n_1},\ldots,x_{n_k})\urcorner)\).  The
sentence coded there is \(\mathsf R_i(c_{n_1},\ldots,c_{n_k})\), so
\(\rho(i,\operatorname{code}_k(\vec n))=1\) iff
\(\vec n\in\mathsf R_i^{\mathcal M}\).

Define \(\zeta(i,z)\) as follows.  If \(i\notin I_F\), or if a decoded entry
lies outside \(M\), return \(0\).  Otherwise search for the least \(m\in M\)
with
\(\ulcorner\mathsf f_i(x_{n_1},\ldots,x_{n_k})=x_m\urcorner\in P\).  The
value \(m=\mathsf f_i^{\mathcal M}(\vec n)\) lies in \(M\) and satisfies the
condition, so the search terminates.  No other \(m\) satisfies it, because
the sentence asserts equality with \(c_m\).  Define \(\gamma(i)=0\) for
\(i\notin I_C\).  Otherwise let \(\gamma(i)\) be the least \(m\in M\) with
\(\ulcorner\mathsf c_i=x_m\urcorner\in P\), which exists and is
\(\mathsf c_i^{\mathcal M}\).  All three functions are total and computable,
since \(I_R,I_F,I_C\), \(M\) and \(P\) are decidable and \(a_R,a_F\) are
computable.

\emph{(c) implies (a).}  For \(n\in\mathrm{Tm}\) with variable indices in
\(M\), define \(\operatorname{val}(n)\) by course-of-values recursion:
\(\operatorname{val}(\ulcorner x_j\urcorner)=j\),
\(\operatorname{val}(\ulcorner\mathsf c_i\urcorner)=\gamma(i)\), and
\[
 \operatorname{val}(\ulcorner\mathsf f_i(t_1,\ldots,t_k)\urcorner)
 =\zeta\bigl(i,\operatorname{code}_k(\operatorname{val}\ulcorner t_1\urcorner,\ldots,\operatorname{val}\ulcorner t_k\urcorner)\bigr).
\]
By induction on terms, \(\operatorname{val}(\ulcorner t\urcorner)\) is the
value in \(\mathcal M_M\) of the closed \(\Lang(M)\)-term obtained from \(t\)
by replacing each \(x_j\) with \(c_j\), and this value lies in \(M\).  The
base cases hold because \(c_j\) denotes \(j\) and
\(\gamma(i)=\mathsf c_i^{\mathcal M}\in M\).  In the step, the arguments lie
in \(M^k\) by the induction hypothesis, and \(\zeta\) agrees with
\(\mathsf f_i^{\mathcal M}\) on \(M^k\), which maps into \(M\).

Now let \(n\) be given.  If \(n\notin\mathrm{Lit}\), or if some variable
index of \(n\) lies outside \(M\), then
\(n\notin\ulcorner\operatorname{Diag}(\mathcal M)\urcorner\).  Both tests are
decidable by \cref{thm:p03-w03-p03-063-parsing} and the decidability of
\(M\).  Otherwise \(\sigma_n\) is defined.  If \(n\) codes \(t=s\),
then \(\sigma_n\) is true iff
\(\operatorname{val}\ulcorner t\urcorner=\operatorname{val}\ulcorner s\urcorner\).
If \(n\) codes \(\mathsf R_i(\vec t)\), then \(\sigma_n\) is true iff
\(\rho(i,\operatorname{code}_k(\operatorname{val}\ulcorner\vec t\urcorner))=1\),
because the values lie in \(M^k\).  A code with tag \(5\) is in the set iff
the atomic code inside it is a false sentence.  This decides
\(\ulcorner\operatorname{Diag}(\mathcal M)\urcorner\).  It is exactly the
procedure of \cref{def:p03-w03-p03-063-decide-literal}.

Condition (a) gives decidability of \(M\) by the first paragraph, and (b) and
(c) include it.

\emph{The binary sequence \(D(\mathcal M)\).}  The set \(\mathrm{At}\) is
decidable.  It is infinite, since the codes \(\ulcorner x_j=x_j\urcorner\) are
pairwise distinct.  So \(i\mapsto a_i\) is computable: search for successive
elements.  Its inverse on \(\mathrm{At}\) is computable: count the elements
of \(\mathrm{At}\) below \(n\).  By definition,
\(D(\mathcal M)(i)=\chi_P(a_i)\).  If (b) holds, \(D(\mathcal M)\) is
therefore computable.  Conversely, if \(D(\mathcal M)\) is computable, then
\(\chi_P(n)\) is \(0\) for \(n\notin\mathrm{At}\) and is \(D(\mathcal M)(i)\)
for \(n=a_i\).  Moreover,
\(j\in M\) iff \(D(\mathcal M)(i)=1\) for the \(i\) with
\(a_i=\ulcorner x_j=x_j\urcorner\).  So (b) holds.
\end{proof}
\end{annexitem}

\begin{annexitem}{annex:p03-w03-p03-063-pullback}{thm:p03-w03-p03-063-pullback}{Transport along computable bijections}
\begin{proof}
For \(m\in\N\), \(g^{-1}(m)\) is the least \(n\) with \(g(n)=m\).  The search
terminates because \(g\) is onto, and it returns the right value because
\(g\) is injective.  So \(g^{-1}\) is computable.

By \cref{thm:p03-w03-p03-063-equivalence}, \(\mathcal M\) has uniformly
computable interpretations \(\rho,\zeta,\gamma\).  For
\(z=\operatorname{code}_k(\vec n)\) with \(k\) the relevant arity, put
\begin{gather*}
 \rho'(i,z)=\rho\bigl(i,\operatorname{code}_k(g\vec n)\bigr),\\
 \gamma'(i)=g^{-1}(\gamma(i)),\\
 \zeta'(i,z)=g^{-1}\bigl(\zeta(i,\operatorname{code}_k(g\vec n))\bigr).
\end{gather*}
Put the value \(0\) at indices outside \(I_R\), \(I_F\) or \(I_C\).  These
functions are total and computable.  By the displayed definition of the
pullback, they are uniformly computable interpretations of
\(g^{-1}(\mathcal M)\), whose domain \(\N\) is decidable.  By
\cref{thm:p03-w03-p03-063-equivalence}, \(g^{-1}(\mathcal M)\) is computable.

The map \(g\) is a bijection.  It satisfies
\(\vec n\in\mathsf R_i^{g^{-1}(\mathcal M)}\iff g\vec n\in\mathsf R_i^{\mathcal M}\),
\(g(\mathsf f_i^{g^{-1}(\mathcal M)}(\vec n))=\mathsf f_i^{\mathcal M}(g\vec n)\)
and \(g(\mathsf c_i^{g^{-1}(\mathcal M)})=\mathsf c_i^{\mathcal M}\).
So it is a surjective embedding, that is, an isomorphism
(\cref{sec:maps-elementarity-01}).

Let \(h:\mathcal A\to\mathcal M\) be a computable isomorphism from an
\(\omega\)-presentation \(\mathcal A\).  Since \(h\) is an isomorphism,
\(\mathsf R_i^{\mathcal A}(\vec n)\) iff \(\mathsf R_i^{\mathcal M}(h\vec n)\),
\(\mathsf f_i^{\mathcal A}(\vec n)=h^{-1}(\mathsf f_i^{\mathcal M}(h\vec n))\)
and \(\mathsf c_i^{\mathcal A}=h^{-1}(\mathsf c_i^{\mathcal M})\).  Thus
\(\mathcal A=h^{-1}(\mathcal M)\), which is computable by the first part.
Finally, a composite of computable total functions is computable, and the
inverse of a computable bijection is computable by the search above.
\end{proof}
\end{annexitem}

\begin{annexitem}{annex:p03-w03-p03-063-domain}{thm:p03-w03-p03-063-domain}{What the domain convention changes}
\begin{proof}
\emph{(a)}  By \cref{thm:p03-w03-p03-063-equivalence}, \(M\) is decidable
and \(\mathcal M\) has uniformly computable interpretations
\(\rho,\zeta,\gamma\).  Let \(h(0)\) be the least element of \(M\), and let
\(h(i+1)\) be the least element of \(M\) above \(h(i)\).  The searches
terminate because \(M\) is infinite, so \(h\) is computable and is a bijection
\(\N\to M\).  On \(M\), the inverse is \(h^{-1}(m)=|\{d<m:d\in M\}|\), a
bounded count, and this formula defines a computable total function on \(\N\).
Transport \(\mathcal M\) along \(h\): put
\(\mathsf R_i^{\mathcal A}(\vec n)\iff\mathsf R_i^{\mathcal M}(h\vec n)\),
\(\mathsf f_i^{\mathcal A}(\vec n)=h^{-1}(\mathsf f_i^{\mathcal M}(h\vec n))\)
and \(\mathsf c_i^{\mathcal A}=h^{-1}(\mathsf c_i^{\mathcal M})\).  As in
\cref{annex:p03-w03-p03-063-pullback}, \(h\) is an isomorphism
\(\mathcal A\to\mathcal M\).  The functions
\(\rho(i,\operatorname{code}_k(h\vec n))\),
\(h^{-1}(\zeta(i,\operatorname{code}_k(h\vec n)))\) and \(h^{-1}(\gamma(i))\)
are uniformly computable interpretations of \(\mathcal A\).  They are correct
because \(\zeta\) and \(\gamma\) take values in \(M\) at arguments from \(M\),
where \(h^{-1}\) is the true inverse.  By
\cref{thm:p03-w03-p03-063-equivalence}, \(\mathcal A\) is computable.

\emph{(b)}  Let \(\mathcal B\) be a finite \(\Lang\)-structure with \(k\)
elements, \(k\geq1\) since domains are nonempty.  List them as
\(b_0,\ldots,b_{k-1}\), and transport \(\mathcal B\) along \(l\mapsto b_l\) to
a structure \(\mathcal M\) with domain \(\{0,\ldots,k-1\}\).  This is a
\((\sqsubseteq\omega)\)-presentation isomorphic to \(\mathcal B\).  Its
domain is finite and so decidable.  Since \(\Lang\) is finite, there are
finitely many symbols, and each interpretation is a finite table on
\(\{0,\ldots,k-1\}\).  Let \(\rho(i,z)\) decode \(z\) to an \(a_R(i)\)-tuple
and look it up in the table of \(\mathsf R_i\) when \(i\in I_R\) and every
entry is below \(k\); otherwise let it return \(0\).  Define \(\zeta\) and
\(\gamma\) in the same way.  These are finite case distinctions over
computable decoding, hence computable.  By
\cref{thm:p03-w03-p03-063-equivalence}, \(\mathcal M\) is computable.

\emph{(c)}  Let \(I_R=\N\), \(a_R(i)=1\) for all \(i\), and
\(I_F=I_C=\varnothing\).  This signature is computable.  Let
\(\mathcal B\) have one element \(b\), with \(\mathsf R_i^{\mathcal B}=\{b\}\)
if \(i\in K\) and \(\mathsf R_i^{\mathcal B}=\varnothing\) otherwise.
Suppose that \(\mathcal M\) is a computable \((\subseteq\omega)\)-presentation
of \(\mathcal B\).  Its domain is \(\{m\}\) for one number \(m\), and
\[
 i\in K\iff \mathsf R_i^{\mathcal M}(m)\iff
 \bigl\langle4,\langle i,\langle0,m\rangle_{\N}\rangle_{\N}\bigr\rangle_{\N}
 \in\ulcorner\operatorname{Diag}(\mathcal M)\urcorner,
\]
using \(\operatorname{code}_1(n)=n\).  For the fixed \(m\), the map
\(i\mapsto\langle4,\langle i,\langle0,m\rangle_{\N}\rangle_{\N}\rangle_{\N}\)
is computable, so \(K\) would be decidable.  This contradicts
\cref{thm:halting-undecidable}.
\end{proof}
\end{annexitem}

\begin{annexitem}{annex:p03-w03-p03-063-no-copy}{thm:p03-w03-p03-063-no-copy}{An abstract countable structure need not have a computable copy}
\begin{proof}
\emph{The structure.}  Take as vertices the pairs \((n,r)\) with \(n\in X\)
and \(r<n+3\).  Let \(E((n,r),(n',r'))\) hold iff \(n=n'\) and
\(r'\equiv r\pm1\pmod{n+3}\).  Since \(n+3\geq3\), each vertex has exactly
two distinct neighbours, \(E\) is symmetric and irreflexive, and the cycle
\(Z_n=\{(n,r):r<n+3\}\) has length \(n+3\).  Since \(X\) is infinite,
\(\mathcal C_X\) is countably infinite.

\emph{A combinatorial claim.}  Suppose that \(w_0,\ldots,w_{n+2}\) are
pairwise distinct vertices of \(\mathcal C_X\) with \(E(w_l,w_{l+1})\) for
\(l<n+2\) and \(E(w_{n+2},w_0)\).  We show that \(n\in X\).  Edges join only
vertices of the same cycle, so every \(w_l\) lies in one cycle \(Z_m\) with
\(m\in X\).  Write its vertices as \(z_q=(m,q)\), indices modulo \(m+3\).
Let \(w_0=z_p\).  Then \(w_1=z_{p+d}\) for some \(d\in\{1,-1\}\).  By
induction, \(w_l=z_{p+ld}\) for \(l\leq n+2\): the vertex \(w_{l+1}\) is a
neighbour of \(z_{p+ld}\), so it is \(z_{p+(l+1)d}\) or
\(z_{p+(l-1)d}=w_{l-1}\), and distinctness excludes \(w_{l-1}\).  The
\(n+3\) distinct vertices lie in a set of \(m+3\) elements, so \(n\leq m\).
The closing edge says that \(z_{p+(n+2)d}\) is adjacent to \(z_p\).  So
\(n+2\equiv\pm1\pmod{m+3}\).  Since \(2\leq n+2\leq m+2\), the case
\(n+2\equiv1\) is impossible.  The case \(n+2\equiv-1\equiv m+2\) gives
\(n=m\in X\).  Conversely, if \(n\in X\), then the vertices of \(Z_n\) in
cyclic order form such a sequence.

\emph{Enumerating \(X\).}  Let \(\mathcal A\) be a computable copy.  By
\cref{thm:p03-w03-p03-063-equivalence}, its domain \(A\) and the relation
\(E^{\mathcal A}\) are decidable.  Let \(Q(n,z)\) say the following: \(z\)
decodes as \(\operatorname{code}_{n+3}(v_0,\ldots,v_{n+2})\); the \(v_l\) lie
in \(A\) and are pairwise distinct; \(E^{\mathcal A}(v_l,v_{l+1})\) holds for
\(l<n+2\); and \(E^{\mathcal A}(v_{n+2},v_0)\) holds.  The relation \(Q\) is
decidable, since it is a finite list of decidable checks after computable
decoding.  An isomorphism preserves and reflects \(E\) and distinctness.  So
by the claim, \(n\in X\) iff \(\exists z\,Q(n,z)\).  The partial function
\(n\mapsto\mu z\,Q(n,z)\) is partial recursive, by closure under unbounded
minimization (\cref{def:p04-partial-recursive}) applied to the computable
characteristic function of \(Q\).  By \cref{thm:p04-tm-partial-recursive}, it
is partial computable.  Its domain is \(X\), so \(X\) is computably
enumerable (\cref{sec:turing-machines-03}).

\emph{The example.}  The set \(\N\setminus K\) is infinite: otherwise \(K\)
would be cofinite, and a cofinite set is decidable by a finite table of
exceptions, contradicting \cref{thm:halting-undecidable}.  By
the corollary in \cref{sec:universal-reductions-03}, \(\N\setminus K\) is not computably enumerable.
So \(\mathcal C_{\N\setminus K}\) has no computable copy.  It is countably
infinite, so it has \(\omega\)-presentations, obtained by transporting it
along any bijection with \(\N\).
\end{proof}
\end{annexitem}

\begin{annexitem}{annex:p03-w03-p03-063-omega-copies}{thm:p03-w03-p03-063-omega-copies}{Computable copies of \((\N,<)\) that are not computably isomorphic}
\begin{proof}
\emph{Step 1: a one-to-one computable listing of \(K\).}  By
\cref{thm:universal-computation}, \(n\mapsto\varphi_n(n)\) is partial
computable, and \(K\) is its domain.  By \cref{thm:p04-tm-partial-recursive},
it is partial recursive.  By \cref{thm:p04-kleene-normal-form}, there are an
index \(e_0\) and a primitive recursive relation \(\mathsf T_1\) such that
\(n\in K\) iff \(\exists y\,\mathsf T_1(e_0,n,y)\).  The relation
\(\mathsf T_1\) is decidable by \cref{sec:turing-machines-02}.  Let
\[
 W=\{\langle n,y\rangle_{\N}:\mathsf T_1(e_0,n,y)\text{ and }
 \neg\mathsf T_1(e_0,n,y')\text{ for all }y'<y\}.
\]
This set is decidable, and it contains exactly one element with first
component \(n\) for each \(n\in K\) and none for \(n\notin K\).  The set \(K\)
is infinite, since a finite set is decidable and \(K\) is not
(\cref{thm:halting-undecidable}).  So \(W\) is infinite.  List \(W\) in
increasing order as \(w_0<w_1<\cdots\), which is computable as in
\cref{annex:p03-w03-p03-063-domain}.  Let \(k_s\) be the first component of
\(w_s\).  Then \(s\mapsto k_s\) is computable and one-to-one, with range
\(K\).

\emph{Step 2: the presentation.}  Put
\(\operatorname{key}(2n)=(n,0)\) and \(\operatorname{key}(2s+1)=(k_s,1)\).
Keys are pairwise distinct, because \(s\mapsto k_s\) is one-to-one.  Let
\(x<_My\) iff \(\operatorname{key}(x)\) precedes \(\operatorname{key}(y)\) in
the lexicographic order of \(\N\times\{0,1\}\).  The lexicographic order is a
strict linear order and \(\operatorname{key}\) is injective, so \(<_M\) is a
strict linear order on \(\N\).  It is decidable, since keys are computable and
compared by finitely many comparisons.  Thus
\(\rho(0,\operatorname{code}_2(x,y))=1\) iff \(x<_My\) defines uniformly
computable interpretations, and \(\mathcal M=(\N,<_M)\) is computable by
\cref{thm:p03-w03-p03-063-equivalence}.

\emph{Step 3: \(\mathcal M\cong(\N,<)\).}  If
\(\operatorname{key}(x)=(n,b)\), every \(<_M\)-predecessor of \(x\) has a key
with first coordinate at most \(n\).  Such elements are the \(n+1\) even
numbers \(2m\) with \(m\leq n\), and the odd numbers \(2s+1\) with
\(k_s\leq n\).  There are at most \(n+1\) of the latter, by injectivity.  So
the set \(\operatorname{pred}(x)\) of predecessors is finite.  Let
\(r(x)=|\operatorname{pred}(x)|\).  If \(x<_My\), then
\(\operatorname{pred}(x)\cup\{x\}\subseteq\operatorname{pred}(y)\) by
transitivity, and \(x\notin\operatorname{pred}(x)\), so \(r(x)<r(y)\).  By
linearity, \(r(x)<r(y)\) iff \(x<_My\), and \(r\) is injective.  The image of
\(r\) is downward closed.  If \(r(x)=p>0\), let \(y\) be the
\(<_M\)-greatest element of the finite nonempty set
\(\operatorname{pred}(x)\).  Then
\(\operatorname{pred}(y)=\operatorname{pred}(x)\setminus\{y\}\), so
\(r(y)=p-1\).  The image is infinite because \(r\) is injective on \(\N\).
An infinite downward-closed subset of \(\N\) is \(\N\).  Hence
\(r:(\N,<_M)\to(\N,<)\) is a bijection that preserves and reflects the order,
which is an isomorphism.

\emph{Step 4: successors.}  The elements strictly between \(2n\) and
\(2n+2\) in \(<_M\) are those with key \((n,1)\).  These are the odd numbers
\(2s+1\) with \(k_s=n\), and there is one exactly when \(n\in K\).  Since
\(2n<_M2n+2\), the immediate \(<_M\)-successor of \(2n\) is \(2n+2\) iff
\(n\notin K\).

\emph{Step 5: no computable isomorphism.}  Suppose that
\(g:(\N,<)\to\mathcal M\) is a computable isomorphism.  By
\cref{thm:p03-w03-p03-063-pullback}, \(g^{-1}\) is computable.  An order
isomorphism carries immediate successors to immediate successors, and the
immediate successor of \(p\) in \((\N,<)\) is \(p+1\).  So the immediate
\(<_M\)-successor of \(x\) is \(g(g^{-1}(x)+1)\).  By Step~4,
\(n\in K\) iff \(g(g^{-1}(2n)+1)\neq2n+2\), which would make \(K\) decidable,
contrary to \cref{thm:halting-undecidable}.  The standard presentation
\((\N,<)\) is computable, since \(x<y\) is decidable.  It is isomorphic to
the computable copy \(\mathcal M\), but not computably.  So \((\N,<)\) is not
computably categorical in the sense of \cref{def:p03-computably-categorical}.

\emph{(b)}  Let \(\pi\) exchange \(2n\) and \(2n+1\) when \(n\in K\), and fix
both when \(n\notin K\).  This \(\pi\) is a bijection of \(\N\).  Put
\(x<'y\) iff \(\pi(x)<\pi(y)\).  Then
\(\pi:(\N,<')\to(\N,<)\) is an isomorphism.  Also \(2n+1<'2n\) iff
\(n\in K\): if \(n\in K\), then \(\pi(2n+1)=2n<2n+1=\pi(2n)\), and otherwise
\(2n+1<2n\) is false.  If \((\N,<')\) were computable, then the computable
map \(n\mapsto\ulcorner c_{2n+1}<c_{2n}\urcorner\), followed by its decision
procedure, would decide \(K\).  So \((\N,<')\) is an \(\omega\)-presentation
of \((\N,<)\) that is not computable.
\end{proof}
\end{annexitem}
